\section{Global Gaussian trace of the original zeta
function}\label{global-gaussian-trace-of-the-original-zeta-function}

24 September 2026. Independent mathematical derivation GGT0--GGT8.

\subsection{GGT0. Source, original objects, and the next
calculation}\label{ggt0.-source-original-objects-and-the-next-calculation}

The preceding \href{GLOBAL_POLYNOMIAL_DUAL_DEFECT.md}{PGD0--PGD7}
construct the actual nonzero class \(c_\zeta=[\lambda_1]\) in the
residue-duality cokernel and its complete real orbit. The next
calculation is a global original-zeta trace that distinguishes the
identity dilation from every other fixed dilation. This note derives
that trace with all prime, endpoint, Gamma, and crossed trivial-zero
terms, then proves the asymptotics needed for the orbit calculation.

The controlling source remains \(Z_0\) absence, primitive \(Z_1/\tau\)
presence without parity or source addition, and the supplied arithmetic
\(Z_2\) data. The source operations, their retractions, and the complete
global corpus arguments referenced in \nolinkurl{CORPUS_AND_OPERATION_RULES.md}
remain controlling. Its later instruction USR-64a88219a3ecddd3 asks what
would advance the calculation using everything already known, and then
asks for that calculation. The constant arithmetic function \(1\) and
the positive real numbers below belong to the already reconstructed
receiving arithmetic. None is identified with primitive \(\tau\).

PGD was read completely for this continuation. The analytic input is the
fully proved original logarithmic-derivative identity
\href{ORIGINAL_ZETA_JACOBIAN_TRACE_BRIDGE.md}{JTB4}, whose complete
proof was independently checked in
\href{ORIGINAL_ZETA_JACOBIAN_TRACE_BRIDGE_INDEPENDENT_CHECK.md}{JTC3}.
It is based on the retained Hadamard product, with every genus-one
exponential and the complete original multiplier. Its human-source
attribution is Jacques Hadamard's factorization, as entered through the
indexed author TeX of Alain Connes,
\href{https://arxiv.org/abs/2602.04022v1}{arXiv:2602.04022v1},
rhready.tex lines526--535. This is the recorded prior source reading,
not a claim to have freshly read the entire paper here.

The original coefficient and analytic contexts remain Alain Connes and
Caterina Consani,
\href{https://arxiv.org/abs/0903.2024v3}{arXiv:0903.2024v3, §5}, and
Ralf Meyer,
\href{https://arxiv.org/abs/math/0412277v3}{arXiv:math/0412277v3}, as
entered in ASD and OMS with their exact source versions. Their geometric
or range results are not being credited with the new Gaussian
calculation. The elementary Euler-product and Gamma estimates required
below are proved explicitly; no complex Stirling estimate with
unspecified branches is assumed.

Write \(\mathscr Z\) for the actual distinct nontrivial zeros of the
original \(\zeta\), with every actual multiplicity \(m_\rho\). Let
\(T\ge1\), \(x>0\), with \(\log x\) the real logarithm. Define the
entire test function and the full zero sum \[
\Phi_{T,x}(s)=x^s e^{s^2/T^2}
=\exp(s\log x+s^2/T^2),
\qquad
S_T(x)=\sum_{\rho\in\mathscr Z}m_\rho x^\rho e^{\rho^2/T^2}.
\tag{GGT0.1}
\] The variable \(T\) is a test-function parameter, not an assigned
clock, coordinate, or weight on primitive \(\tau\).

\subsection{GGT1. Convergence and the full original logarithmic
derivative}\label{ggt1.-convergence-and-the-full-original-logarithmic-derivative}

The original receiving test algebra is \[
\mathcal B=\{F\text{ entire}:
b_{A,N}(F)=\sup_{|\sigma|\le A,t\in\mathbb R}
(1+|t|)^N|F(\sigma+it)|<\infty
\text{ for every }A,N\}.
\tag{GGT1.1}
\] For fixed \(T,x\), the exact modulus is \[
|\Phi_{T,x}(\sigma+it)|
=x^\sigma e^{\sigma^2/T^2}e^{-t^2/T^2},
\tag{GGT1.2}
\] so \(\Phi_{T,x}\in\mathcal B\). The already proved zero count
\(n(R)\le C(R+2)^{3/2}\), with multiplicities, gives absolute
convergence of (GGT0.1), since \[
|x^\rho e^{\rho^2/T^2}|
\le\max(1,x)e^{1/T^2}e^{-(\Im\rho)^2/T^2}.
\tag{GGT1.3}
\] Summation over integer height intervals proves the claim from the
polynomial count and Gaussian decay. In particular no ordering of
conditionally convergent zero terms is being chosen.

For an integrable meromorphic function on the two edges, retain \[
\mathfrak C(F)=\frac1{2\pi i}
\left(\int_{2-i\infty}^{2+i\infty}F(s)\,ds
-\int_{-1-i\infty}^{-1+i\infty}F(s)\,ds\right).
\tag{GGT1.4}
\] Both integrals run upward. JTB4 proves, for every
\(\Phi\in\mathcal B\), \[
\mathfrak C\!\left(\Phi\,\frac{\zeta'}{\zeta}\right)
=\sum_{\rho\in\mathscr Z}m_\rho\Phi(\rho)-\Phi(1).
\tag{GGT1.5}
\] For clarity, its exact retained factor comparison is \[
F_*(s)=\frac{s(s-1)}8\pi^{-s/2}\Gamma(s/2)\zeta(s),
\] \[
\frac{F_*'}{F_*}
=b_0+\sum_{\rho\in\mathscr Z}m_\rho
\left(\frac1{s-\rho}+\frac1\rho\right),
\qquad b_0=\tfrac12\log(4\pi)-1-\tfrac\gamma2,
\] \[
\frac{F_*'}{F_*}
=\frac{\zeta'}{\zeta}+\frac1s+\frac1{s-1}
-\frac12\log\pi+\frac12\psi(s/2),
\qquad \psi=\Gamma'/\Gamma.
\tag{GGT1.6}
\] On the two edges the sum of absolute values of the paired genus-one
terms is \(O((1+|t|)^{3/2})\), by the retained zero count and
\(|s-\rho|\ge1\). Thus it can be integrated term by term against
\(\Phi\). Each paired term gives its single residue, and the multiplier
derivative in the last line has cancelling residues at zero and residue
\(+1\) at one. Its subtraction gives the original pole contribution
\(-\Phi(1)\) in (GGT1.5). These are the actual reasons that both the
zero sum and the edge integrals are absolutely convergent. \(F_*\) has
not replaced \(\zeta\) in (GGT1.5).

Applied to (GGT0.1), the exact starting equation is \[
S_T(x)=x e^{1/T^2}
+\mathfrak C\!\left(\Phi_{T,x}\,\frac{\zeta'}{\zeta}\right).
\tag{GGT1.7}
\]

\subsection{GGT2. The Gaussian Mellin integral and the original prime
series}\label{ggt2.-the-gaussian-mellin-integral-and-the-original-prime-series}

For every real \(c\) and \(y>0\), direct parameterization gives \[
\begin{aligned}
\frac1{2\pi i}\int_{c-i\infty}^{c+i\infty}
y^s e^{s^2/T^2}\,ds
&=\frac{y^ce^{c^2/T^2}}{2\pi}
\int_{\mathbb R}e^{-t^2/T^2}
e^{it(\log y+2c/T^2)}\,dt\\
&=\frac{T}{2\sqrt\pi}\,
y^ce^{c^2/T^2}
\exp\left[-\frac{T^2}{4}
\left(\log y+\frac{2c}{T^2}\right)^2\right]\\
&=\frac{T}{2\sqrt\pi}
\exp\left[-\frac{T^2}{4}\log^2y\right].
\end{aligned}
\tag{GGT2.1}
\] The last equality retains and evaluates both cross terms; no factor
depending on \(c\) or \(y\) has been omitted. The Fourier Gaussian used
in the second line can be verified by differentiating its integral in
the frequency \(\kappa\). Integration by parts gives
\(J'(\kappa)=-T^2\kappa J(\kappa)/2\), while the Gaussian integral gives
\(J(0)=T\sqrt\pi\). Solving this scalar differential equation proves
exactly the displayed factor and exponent.

Define the original von Mangoldt coefficients by \[
\Lambda(n)=
\begin{cases}\log p,&n=p^k,\ p\text{ prime},\ k\ge1,\\0,&\text{otherwise}.\end{cases}
\tag{GGT2.2}
\] They are defined using the retained original arithmetic primes, not a
substituted spectrum. On \(\Re s>1\), the absolutely convergent Euler
product yields \[
-\frac{\zeta'(s)}{\zeta(s)}
=\sum_p\sum_{k\ge1}(\log p)p^{-ks}
=\sum_{n\ge2}\Lambda(n)n^{-s}.
\tag{GGT2.3}
\] For the differentiation, restrict first to any
\(\Re s\ge1+\varepsilon\). The logarithmic series
\(\sum_{p,k}p^{-ks}/k\) and its differentiated series converge uniformly
there, since the latter is bounded by
\(\sum_{n\ge2}(\log n)n^{-1-\varepsilon}\). This proves the identity
before evaluating either edge.

The original functional equation, with every Gamma factor, is \[
\zeta(s)=\chi(s)\zeta(1-s),\qquad
\chi(s)=\pi^{s-1/2}
\frac{\Gamma((1-s)/2)}{\Gamma(s/2)}
=2^s\pi^{s-1}\sin(\pi s/2)\Gamma(1-s).
\tag{GGT2.4}
\] Its logarithmic derivative is \[
\frac{\zeta'(s)}{\zeta(s)}
=g(s)-\frac{\zeta'(1-s)}{\zeta(1-s)},
\qquad
g(s)=\frac{\chi'(s)}{\chi(s)}
=\log\pi-\frac12\psi((1-s)/2)-\frac12\psi(s/2).
\tag{GGT2.5}
\] On the left edge \(\Re s=-1\), equations (GGT2.3) and (GGT2.5) give
\[
\frac{\zeta'(s)}{\zeta(s)}
=g(s)+\sum_{n\ge2}\frac{\Lambda(n)}n n^s.
\tag{GGT2.6}
\] Each series can be interchanged with its edge integral: on the right
its absolute coefficient sum is \(\sum\Lambda(n)n^{-2}<\infty\); on the
left the coefficient absolute values are again \(\Lambda(n)n^{-2}\). The
remaining test function is integrable by (GGT1.2). This proves an
absolute Fubini bound before the Gaussian integral is evaluated.

\subsection{GGT3. The exact trace, with its full archimedean
term}\label{ggt3.-the-exact-trace-with-its-full-archimedean-term}

Define, on the original left edge, \[
A_T(x)=\frac1{2\pi i}
\int_{-1-i\infty}^{-1+i\infty}
x^s e^{s^2/T^2}\frac{\chi'(s)}{\chi(s)}\,ds .
\tag{GGT3.1}
\] The estimates below prove absolute convergence directly. Substituting
(GGT2.3), (GGT2.6), and (GGT2.1) into the two distinct integrals of
(GGT1.7) proves \[
\boxed{
\begin{aligned}
S_T(x)= {}&x e^{1/T^2}\\
&-\frac{T}{2\sqrt\pi}
\sum_{n\ge2}\Lambda(n)\left[
\exp\left(-\frac{T^2}{4}\log^2(x/n)\right)
+\frac1n\exp\left(-\frac{T^2}{4}\log^2(xn)\right)
\right]\\
&-A_T(x).
\end{aligned}}
\tag{GGT3.2}
\] The sign of the second prime family is negative because it is a
positive series on the subtracted left edge. Its coefficient is exactly
\(1/n\). The original pole term is exactly \(x e^{1/T^2}\).

Both prime sums converge absolutely. In addition to the Fubini argument,
their final terms are bounded by constant multiples of
\((\log n)\exp[-c(\log n)^2]\) for fixed \(x,T\); this is summable, for
example by comparison with \(n^{-2}\) after increasing a finite cutoff.
Every prime and prime-power coefficient remains.

The original endpoint and trivial-zero information can also be exposed
by finite contour shifts, without replacing (GGT3.1). The meromorphic
function \(g=\chi'/\chi\) has residue \(+1\) at \[
0,-2,-4,\ldots
\] and residue \(-1\) at \[
1,3,5,\ldots.
\] These follow from the simple Gamma poles in the numerator and
denominator of the full quotient (GGT2.4). Define \(A_T^{(c)}(x)\) by
the same integral as (GGT3.1) on any vertical edge \(c\) avoiding those
real poles. For a fixed finite integer \(N\ge0\), the residue theorem
gives \[
A_T(x)
=A_T^{(-2N-1)}(x)
+\sum_{r=1}^{N}x^{-2r}e^{4r^2/T^2},
\tag{GGT3.3}
\] \[
A_T(x)=A_T^{(1/2)}(x)-1.
\tag{GGT3.4}
\] The second identity crosses the pole of \(g\) at zero with residue
\(\Phi_{T,x}(0)=1\). Thus the equivalent identity \[
S_T(x)=1+x e^{1/T^2}
-\text{the two full prime sums in (GGT3.2)}
-A_T^{(1/2)}(x)
\tag{GGT3.5}
\] has both endpoint contributions displayed. Equations
(GGT3.3)--(GGT3.4) are exact comparisons, not discarded terms. The
finite-shift horizontal integrals tend to zero by Gaussian decay and the
fixed-strip Gamma estimates proved in GGT4.

The negative even points in (GGT3.3) are the original simple trivial
zeros: (GGT2.4), together with \(\zeta(1+2r)\ne0\), gives \[
\zeta'(-2r)=(-1)^r2^{-2r-1}\pi^{-2r}(2r)!\zeta(1+2r)\ne0.
\tag{GGT3.6}
\] Their logarithmic-derivative residue is one, explaining the
coefficient in (GGT3.3). An infinite left shift is not justified: for
fixed \(T\), the terms \(x^{-2r}e^{4r^2/T^2}\) eventually grow. The full
remaining contour in every finite identity is retained.

\subsection{GGT4. A complete branch-controlled Gamma
estimate}\label{ggt4.-a-complete-branch-controlled-gamma-estimate}

We prove constants \(C_0,C_1,C_2\) such that, for \(t\in\mathbb R\), \[
|g(-1+it)|\le C_0\log(2+|t|),\qquad
\left|\frac d{dt}g(-1+it)\right|
\le\frac{C_1}{1+|t|},
\tag{GGT4.1}
\] and, for \(|t|\ge2\), \[
g(-1+it)=\log(2\pi)-\log|t|+r(t),
\qquad |r(t)|\le C_2/|t|.
\tag{GGT4.2}
\] The remainder can be complex; the formula does not silently replace
the original function by its real part.

Here is an elementary integral estimate for \(\psi\). For \(\Re z>0\),
use the principal logarithm, whose argument lies in \((-\pi/2,\pi/2)\).
Euler's Gamma limit gives \[
\psi(z)=\lim_{N\to\infty}
\left(\log N-\sum_{n=0}^{N-1}\frac1{n+z}\right).
\tag{GGT4.3}
\] The limit follows, with locally uniform first derivatives, from \[
\Gamma(z)=\lim_{N\to\infty}
\frac{N!\,N^z}{z(z+1)\cdots(z+N)}.
\] For completeness, integration by parts in the beta integral gives the
displayed finite ratio before the factor \(N^z\). After the substitution
\(v=Nu\), that factor times the beta integral is
\(\int_0^N v^{z-1}(1-v/N)^N\,dv\). The bound \((1-v/N)^N\le e^{-v}\) on
\(0\le v<N\), including logarithmic factors for derivatives, proves
locally uniform convergence to the defining Gamma integral. Its
logarithmic derivative gives (GGT4.3); the extra final summand with
index \(N\) tends to zero.

Let \(\widetilde B_1(u)=u-\lfloor u\rfloor-\tfrac12\) off the integers.
It is bounded in absolute value by \(1/2\). On each integer interval,
integrate \(f(u)=(u+z)^{-1}\) once against this function. Summing the
endpoint terms yields \[
\sum_{n=0}^{N-1}\frac1{n+z}
=\Log(N+z)-\Log z+\frac1{2z}-\frac1{2(N+z)}
-\int_0^N\frac{\widetilde B_1(u)}{(u+z)^2}\,du.
\] All paths \(u+z\) lie in the right half-plane, so the logarithms here
use the same fixed branch. Taking the limit in (GGT4.3) proves \[
\psi(z)=\Log z-\frac1{2z}
+\int_0^\infty\frac{\widetilde B_1(u)}{(u+z)^2}\,du,
\] \[
\psi'(z)=\frac1z+\frac1{2z^2}
-2\int_0^\infty\frac{\widetilde B_1(u)}{(u+z)^3}\,du.
\tag{GGT4.4}
\] Absolute convergence on every compact subset of \(\Re z>0\) justifies
differentiation.

If \(z=c+iy\), \(c>0\), \(|y|\ge1\), then \[
\int_0^\infty\frac{du}{|u+z|^2}
\le\frac{\pi}{2|y|},\qquad
\int_0^\infty\frac{du}{|u+z|^3}
\le\frac1{|y|^2}.
\] Consequently (GGT4.4) gives \[
\psi(c+iy)=\Log(c+iy)+O(|y|^{-1}),\qquad
\psi'(c+iy)=\frac1{c+iy}+O(|y|^{-2}),
\tag{GGT4.5}
\] uniformly for \(c\) in any fixed compact interval of positive real
numbers. These estimates have been proved from the integral, rather than
taken from an unspecified angular sector in a Stirling formula.

For the actual line \(s=-1+it\), put \[
y=t/2,\qquad z_1=1-iy,\quad z_2=\tfrac12+iy,\quad z_0=-\tfrac12+iy.
\] The recurrence \(\psi(z+1)=\psi(z)+1/z\) and (GGT2.5) give the exact
identity \[
g(-1+it)=\log\pi-\frac12\psi(z_1)
-\frac12\psi(z_2)+\frac1{2z_0}.
\tag{GGT4.6}
\] Both digamma arguments in this expression have positive real part,
and \(z_0\) is never zero for real \(t\).

For \(|y|\ge1\), the principal logarithms satisfy \[
\Log z_1+\Log z_2
=2\log|y|
+\Log\left(1+\frac1{2y^2}+\frac{i}{2y}\right).
\tag{GGT4.7}
\] Indeed \(z_1z_2=y^2+\tfrac12+\tfrac{i}{2}y\) has positive real part,
and the sum of their arguments is the principal argument of that
product. No multiple of \(2\pi i\) enters. Along the segment from \(1\)
to \(1+w\), with \(w=1/(2y^2)+i/(2y)\), the real part is at least one.
Hence the last logarithm has absolute value at most \(|w|\le |y|^{-1}\).
Equations (GGT4.5)--(GGT4.7) prove (GGT4.2), including both signs of
\(t\).

Differentiating the exact (GGT4.6) gives \[
\frac d{dt}g(-1+it)
=\frac i4\psi'(z_1)-\frac i4\psi'(z_2)
-\frac i{4z_0^2}.
\tag{GGT4.8}
\] Equation (GGT4.5) bounds this by \(C/|t|\) for \(|t|\ge2\). On the
compact interval \([-2,2]\), every term and its derivative is
continuous, so it is bounded. This proves (GGT4.1) globally. The same
recurrence argument shifts both digamma arguments into a fixed
positive-real-part strip for any other fixed bounded range of \(\Re s\),
at large \(|\Im s|\). It proves \(g(s)=O(\log(2+|\Im s|))\) uniformly on
that range, which justifies the finite shifts in GGT3.

\subsection{GGT5. The archimedean integral away from the
identity}\label{ggt5.-the-archimedean-integral-away-from-the-identity}

Parameterize the exact left-edge integral without discarding its phase:
\[
A_T(x)=\frac{x^{-1}e^{1/T^2}}{2\pi}
\int_{\mathbb R}e^{-t^2/T^2}
e^{it(\log x-2/T^2)}g(-1+it)\,dt.
\tag{GGT5.1}
\] GGT4 and Gaussian decay prove absolute convergence. Let \[
h_T(t)=e^{-t^2/T^2}g(-1+it),\qquad
\kappa_T=\log x-2/T^2.
\] Its derivative is integrable, and (GGT4.1) gives \[
\begin{aligned}
\int_{\mathbb R}|h_T'(t)|\,dt
&\le
C\int_{\mathbb R}\frac{2|t|}{T^2}
e^{-t^2/T^2}\log(2+|t|)\,dt\\
&\quad+C\int_{\mathbb R}\frac{e^{-t^2/T^2}}{1+|t|}\,dt\\
&\le C'\log(2+T).
\end{aligned}
\tag{GGT5.2}
\] For the first bound substitute \(t=Tu\) and use
\(\log(2+T|u|)\le\log(2+T)+\log(1+|u|)\) for \(T\ge1\). For the second,
split at \(|t|=T\); the inner integral is bounded by \(2\log(1+T)\), and
substitution in the outer Gaussian tail gives a constant. This proves
the last line with constants independent of \(T\).

For fixed \(x\ne1\), all sufficiently large \(T\) satisfy \[
|\kappa_T|\ge|\log x|/2.
\] The boundary values of \(h_T\) vanish. Integration by parts in
(GGT5.1), with the exact frequency \(\kappa_T\), therefore gives \[
\left|\int_{\mathbb R}h_T(t)e^{i\kappa_Tt}\,dt\right|
\le\frac{2C'}{|\log x|}\log(2+T).
\] For the remaining bounded interval of \(T\ge1\), the original
absolute integral in (GGT5.1) is uniformly bounded by a fixed integrable
Gaussian times \(\log(2+|t|)\). Increasing the constant thus proves \[
\boxed{A_T(x)=O_x(\log(2+T))\quad(x\ne1,\ T\ge1).}
\tag{GGT5.3}
\] In particular any finite value of \(T\) at which \(\kappa_T=0\) has
been treated by the absolute integral, not by an invalid division by
zero.

\subsection{GGT6. The archimedean integral at the identity, with its
next
constant}\label{ggt6.-the-archimedean-integral-at-the-identity-with-its-next-constant}

At \(x=1\), write for \(t\ne0\) \[
g(-1+it)=\log(2\pi)-\log|t|+r(t).
\] By GGT4, \(|r(t)|\le C/|t|\) for \(|t|\ge2\). For \(|t|<2\),
continuity of \(g\) gives an integrable bound \(C+|\log|t||\).
Consequently \[
\int_{\mathbb R}e^{-t^2/T^2}|r(t)|\,dt=O(\log(2+T)).
\tag{GGT6.1}
\] The exact phase in (GGT5.1) obeys \[
|e^{-2it/T^2}-1|\le2|t|/T^2.
\] After multiplication by the Gaussian and \(1+|\log|t||\), its
integral is \(O(\log(2+T))\), by \(t=Tu\). The prefactor
\(e^{1/T^2}=1+O(T^{-2})\) changes an \(O(T\log(2+T))\) integral by
\(O(\log(2+T)/T)\). Thus every error has been controlled before
evaluating the leading integral, and \[
A_T(1)=\frac1{2\pi}
\int_{\mathbb R}e^{-t^2/T^2}
\bigl(\log(2\pi)-\log|t|\bigr)\,dt
+O(\log(2+T)).
\tag{GGT6.2}
\] The two exact Gaussian integrals are \[
\int_{\mathbb R}e^{-t^2/T^2}\,dt=T\sqrt\pi,
\] \[
\int_{\mathbb R}e^{-t^2/T^2}\log|t|\,dt
=T\sqrt\pi\log T+\frac T2\Gamma'(1/2).
\tag{GGT6.3}
\] For the second, substitute \(t=Tu\), then \(v=u^2\) on the two
half-lines. Differentiation of the defining Gamma integral at \(1/2\) is
justified by the integrable logarithmic factor.

The Euler limit (GGT4.3) at \(z=1/2\) gives \[
\psi(1/2)=
\lim_{N\to\infty}\bigl(\log N-2H_{2N}+H_N\bigr)
=-\gamma-2\log2.
\tag{GGT6.4}
\] Here \(H_N=\sum_{n=1}^N1/n\) and \(\gamma=\lim(H_N-\log N)\). Since
\(\Gamma(1/2)=\sqrt\pi\), substituting (GGT6.3)--(GGT6.4) proves \[
\boxed{
A_T(1)=\frac{T}{2\sqrt\pi}
\left[-\log T+\log(4\pi)+\frac\gamma2\right]
+O(\log(2+T)).
}
\tag{GGT6.5}
\] The original phase and full Gamma terms have been bounded, not
omitted without a comparison.

\subsection{GGT7. All prime resonances and the global trace
bounds}\label{ggt7.-all-prime-resonances-and-the-global-trace-bounds}

For fixed \(x>0\), the prime sum in (GGT3.2) has the \(T\)-independent
summable majorant \[
\sum_{n\ge2}(\log n)
\left[e^{-\log^2(x/n)/4}
+n^{-1}e^{-\log^2(xn)/4}\right]<\infty,
\quad T\ge1.
\tag{GGT7.1}
\] This proves immediately that its full contribution, including the
factor \(T/(2\sqrt\pi)\), is \(O_x(T)\). Combined with GGT5 and the
bounded original pole term, it gives \[
\boxed{S_T(x)=O_x(T)\quad(x\ne1).}
\tag{GGT7.2}
\]

There is a sharper statement retaining each possible exact resonance.
Define \[
R(x)=
\begin{cases}\Lambda(x),&x\text{ is an integer }\ge2,\\0,&\text{otherwise}\end{cases}
\ +\
\begin{cases}x\Lambda(1/x),&1/x\text{ is an integer }\ge2,\\0,&\text{otherwise}.\end{cases}
\tag{GGT7.3}
\] The two alternatives cannot both be nonzero. These are exactly the
terms with \(\log(x/n)=0\) or \(\log(xn)=0\) in the two prime families.

All other logarithms in those two families have a positive minimum
absolute value \(\delta_x>0\). Indeed each family tends in absolute
value to infinity with \(n\), and only finitely many indices lie in any
bounded interval; the zero logarithms have already been excluded. For a
nonzero such logarithm \(u\), \(T\ge1\), and \(|u|\ge\delta_x\), \[
\frac{T^2u^2}{4}
\ge\frac{T^2\delta_x^2}{8}+\frac{u^2}{8}.
\] The corresponding series with \(e^{-u^2/8}\) is summable by the same
argument as (GGT7.1). Therefore the full prime sum, before its
prefactor, is \[
R(x)+O_x(e^{-T^2\delta_x^2/8}).
\tag{GGT7.4}
\] For \(x\ne1\), equations (GGT3.2), (GGT5.3), and (GGT7.4) prove the
stronger expansion \[
\boxed{
S_T(x)=-\frac{T}{2\sqrt\pi}R(x)+x
+O_x(\log(2+T)).
}
\tag{GGT7.5}
\] Thus nontrivial prime-power resonances have been retained rather than
put into an unspecified error term.

At \(x=1\) there is no zero logarithm, and \(\delta_1=\log2\). The prime
contribution is \[
O\!\left(T e^{-T^2(\log2)^2/8}\right).
\tag{GGT7.6}
\] Combining this with \(e^{1/T^2}=1+O(T^{-2})\) and the full constant
in (GGT6.5) gives \[
\boxed{
S_T(1)=
\frac{T}{2\sqrt\pi}
\left[\log T-\log(4\pi)-\frac\gamma2\right]
+1+O(\log(2+T)).
}
\tag{GGT7.7}
\] In particular \[
S_T(1)=\frac{T}{2\sqrt\pi}\log T+O(T).
\tag{GGT7.8}
\] These estimates use the actual zero multiset and the original
arithmetic prime series. They make no assumption about the real parts of
the zeros beyond the established critical-strip inclusion.

\subsection{GGT8. Exact trace tests for the original multiplier
comparison}\label{ggt8.-exact-trace-tests-for-the-original-multiplier-comparison}

The estimates apply to the actual multiplier domain in PGD. If
\(F\in\mathcal B\), then for each fixed \(x>0\), \[
\sup_{T\ge1}
\sum_{\rho\in\mathscr Z}m_\rho
|x^\rho F(\rho)e^{\rho^2/T^2}|<\infty.
\tag{GGT8.1}
\] Indeed \(|x^\rho|\le\max(1,x)\), \(|e^{\rho^2/T^2}|\le e\), and the
rapid bound for \(F(\rho)\), together with the polynomial zero count,
makes the remaining sum absolutely convergent. Thus an actual test
representative has uniformly bounded trace under these Gaussian tests,
while (GGT7.7) identifies the distinct \(T\log T\) behavior of the
arithmetic identity multiplier.

This supplies the global estimate required to test a finite relation
among PGD's already constructed real-orbit multipliers. It does not
replace the original quotient by a sequence space, identify primitive
\(\tau\) with a constant arithmetic function, or assert a purity bound.
The exact prime, Gamma, endpoint, and finite trivial-zero comparisons
are (GGT3.2)--(GGT3.6); the inequalities follow from those full
formulas.
